\renewcommand{\thechapter}{1}
\chapter{Introduction}
\label{chapter:10-Introduction}

%Graph Problem Definition 
\nrscomment{cut this chunk, it's vague and confusing as the entry point}
When solving problems, it is wise to consider the context of the evidence. 
At the micro-level, relationships between pieces of evidence can indicate similarity, dependence or some other link worthy of exploration. 
At the macro level, the evidence might be part of a particular cluster or some other structural feature that provides important context.

Graphs are powerful data structures for representing contextual information and solving problems. 
They explicitly encode the evidence as vertices (nodes) and the relationships as edges (arcs, links).
Intelligent algorithms can solve many problems on graphs but are limited in the size of the graphs they can process. 
In the real world, graphs quickly reach billions or trillions of vertices, \nrscomment{"and Graph-Based Artificial Intelligence, or \gai algorithms are frequently quadratic or worse in complexity."} and the complexity of intelligent algorithms that operate on graphs (Graph-Based Artificial Intelligence, or \gai) is frequently quadratic or worse. 
In addition to the algorithmic complexity, as graphs become larger and take on more complicated shapes, modern general-purpose processors become memory-bound.
\nrscomment{The transition of the sentence gives the expectation that what will follow will be of the form "there is also memory complexity" but you have it flipped with "processors become memory bound". Can you change either the beginning or the end of the sentence to make them cohesive? Something like "Further, modern general purpose processors become memory bound as graphs become larger and take on more complicated shapes." would even do.}

\nrscomment{the content here is good but it reads like a collection of sentences and i have to work hard to track them all and figure out what youre trying to say. try rewriting it}
\gai slows down as graphs get large and complex because of how graphs are stored in memory. 
Most algorithms exploit the graph's structure, and the processing follows the graph's edges. 
Graphs are typically stored in memory in the order vertices were created, so topologically close vertices are not necessarily proximal in memory. 
As a result, prefetching and caching techniques that exploit locality in general-purpose computing are ineffective in graph processing.
\gai must retrieve data required for the next step of the algorithm from the main memory rather than the cache, so the latency and bandwidth of those reads become the bottleneck. 
\nrscomment{this next sentence doesnt really connect directly to the previous paragraph in an obvious way- make it more explicit how this has to do with what you just said}
Real-world application of \gai requires alleviating the bottlenecks in large, complicated graphs. 

Most competing approaches to speed up graph processing focus on generating improvement for a subset of problems and do not necessarily yield general improvement in all \gai performance. 
\nrscomment{next 2-3 sentences sound like speech, they depend on each other to make sense grammatically (i.e. sentence 2 cant stand alone). clean them up}
Accelerator-based approaches like distributed processing, upgrading hardware and adding specialized accelerators provide one path. 
Representation-based approaches like compression, custom data structures and alternate abstraction models are a second.
Algorithmic approaches provide a third path that seeks to exploit parallelism or develop approximation methods to reduce computation time. 
These approaches are interdependent, and decisions about optimizing hardware can limit choices about parallelism or programming models, making it difficult to determine the optimal configuration for a given collection of problems. 

\nrscomment{change the first person ref to the engineers/etc. that are your users}
This thesis explores how to compare each improvement technique's impact on a system that solves real-world problems with \gai so that we can consistently improve \gai performance while justifying our design decisions empirically. 
For example, determining whether adding hardware accelerators like \ac{GPU} or distributing computing across the existing infrastructure is more cost-effective to reduce the processing time. \nrscomment{that wasnt a sentence, theres no subject}
Further, developing a repeatable process to collect the data needed to make these decisions yields an extensible framework to evaluate the full range of \gai problems.
\nrscomment{that last chunk was your contributions summary, need to make it stronger. Use words like 'contributes', 'develops', 'proposes', 'provides' instead of 'explores'. Second and third sentences should be stronger too- "Our framework enables..." and that sort of wording, not "developing x yields y"-- that doesnt neccessarily imply that this thesis develops it.}

\section{Introduction to the Problem}
\label{section:11-introduction-to-problem}
\nrscomment{i think this is the first time you use the word benchmark- without defining or introducing it at all. probably need to work it into the earlier section}
The \nrscomment{Our} \gai Benchmarking Framework provides a mechanism to determine \textit{if}, \textit{how} and \textit{when} to speed up problem-solving on linked data given a set of operating constraints. 
It builds on prior work in benchmarking and workload analyses to create a balanced benchmark that is problem-centric and accelerator-agnostic and evaluates both the \textit{efficiency} and \textit{effectiveness} of a system under test, bringing immediate benefit to analysts, engineers and researchers. 

\subsection{Problem Definition}
This thesis aims to determine \textit{if}, \textit{how} and \textit{when} to speed up problem-solving on linked data given a set of operating \textit{constraints}.
\nrscomment{next 3 sentences can probably be combined into 1-2 and worded to avoid "the how dimension explores how to" etc.}
The \textit{how} dimension explores how to effectively combine the representations, hardware accelerators and algorithms to speed up \gai.
The \textit{if} and \textit{when} dimensions answer when it is appropriate to exploit acceleration to speed up the time to solution. 
In the real world, \textit{constraints} on the efficiency, effectiveness and cost of a solution force a user of \gai \nrscomment{missing a word "to"?} directly impact the if, how and when dimensions, and so are explicitly considered here.

\subsection{Prior Work}
Prior approaches to identifying if, how, and when to accelerate a \gai system \nrscomment{"do not balance the utility..." onward is unclear. Prior approaches dont balance the utility to use for benchmarking of \gai scale??} do not balance the general utility and depth of analysis to use for benchmarking of \gai at scale. 
\textit{Benchmarking} is a process by which systems or components are compared according to some criteria, which is used to determine their relative performance~\cite{Kistowski2015}. \nrscoment{the defn comes too late}
Most benchmarks for \gai take a `black-box' approach to competing systems based on metrics of \textit{efficiency} or \textit{effectiveness} but rarely both, and never their interactions.
They have typically focused on evaluating isolated \textit{kernel} level tasks rather than more complicated problems consistent with the real world. 
They typically use a small selection of datasets that do not necessarily represent the real world and do not cover the breadth of expected input sizes and shapes. 
More support is needed for dynamic graph processing, and no work has been done comparing dynamic and static approaches under the same evaluation framework. 

\nrscomment{i think this is first time you mention workload analysis- need to transition it better "Alternatively, workload..."}
Workload Analyses are `white-box' approaches that analyze the internal state of a system while completing a workload to identify bottlenecks. 
While they yield detailed information, they require access to the system under test for profiling and access to the code to add instrumentation. 
The invasive configuration requirements mean that \nrscomment{drop the definite article} the workload analyses cannot be scaled for general-purpose comparisons and are unsuitable for the broad coverage the \gai benchmark framework pursues. \nrscomment{anthropomorphic- a benchmark cannot pursue broad coverage}

\subsection{The \gai Benchmark Framework}
We introduce the \gai Benchmark Framework as a solution that balances the generality of competitive benchmarks and the detail of workload analyses without requiring invasive system access. 
We leverage constraint theory in conjunction with an empirical approach to algorithmic analysis to develop inputs that specifically trigger bottleneck behavior. 
We then use our observations of those behaviors to develop a `grey-box' understanding of the system's internal state. 
Our approach is \textit{problem-centric}, meaning that the benchmark tasks are real-world problem statements rather than simplified kernels. 
The problem-centric view encourages innovative approaches by not constraining the benchmark to specific approaches. 
It also better captures the deployment of a system in the real world. 
The framework adopts an Input/Output that maximizes the degree of freedom of a benchmark user who can vary their hardware, software stack and implementation without having to change how their code interfaces with the benchmark, allowing for easy experimentation about different combinations of acceleration techniques to determine what works best for their use case. 
Finally, our evaluation approach explores the interaction between efficiency and effectiveness, recognizing that \nrscomment{a fast partial solution can be more useful than a slow complete one.} sometimes a fifty percent solution on time is better than a hundred percent solution too late. 
The evaluation approach also makes it possible to evaluate dynamic and static approaches using the same framework. \nrscomment{dont use 'evaluate' and 'approach' both twice in the same sentence}

\osullikomment{Add example here for problem vs kernel}

\subsection{Impact}
The \gai Benchmark Framework is a flexible tool for experimentation and exploratory analysis for analysts, engineers, and researchers.
Analysts who routinely work with graph data and apply \gai to solve problems can \nrscomment{use the \gai benchmark to} better understand the sizes and shapes of inputs that will impact their work performance and use prior results to build estimates for \gai execution time. 
Engineers can use the \gai benchmark to evaluate which systems they should procure and \nrscomment{to} make architectural decisions about how their system implements \gai to maximize the efficiency and effectiveness across the breadth of tasks the system must support.
Researchers building \gai can use the \gai Benchmark Framework to experimentally vary the hardware and implementation details to find the bottlenecks in \gai as a function of the input datasets and real-world problems it solves. 
The ability to satisfy the needs of these three distinct user groups is a crucial feature of the \gai benchmark framework. \nrscomment{flip the sentence "It is critical that..."}



\section{Our Graph-Based Artificial Intelligence Taxonomy}
\label{section:12-graph-based-ai}
The \gai Benchmark Framework defines Graph-Based Artificial intelligence (\gai) as a taxonomy of intelligence algorithms and machine learning models that learn and reason over information represented in graph data structures.

    
    \begin{figure}[h!]
        \begin{center}
            \includesvg[width=\columnwidth]{SVGs/gai_taxonomy}
            \caption{The \gai Taxonomy with example tasks}
            \label{figure:gai-taxonomy}
        \end{center}
    \end{figure}



Prior efforts to define a taxonomy of \gai focus primarily on \textit{Graph Neural Network}~\cite{Smith2021, Zhou2020a}, \textit{Graph Mining}~\cite{Chen2022a} or specific problems like signal processing~\cite{Xia2021}.
We extend these taxonomies to include:
\begin{itemize}
    \item \textbf{Graph Problem Solving \ac{GPS}.} Classic AI tasks that formulate the world as a graph, including navigation and critical path analysis. 
    \item \textbf{Graph Mining (\ac{GM}).} Analyze the structural characteristics of a graph, including spatial pattern matching and community detection. 
    \item \textbf{Graph Reasoning (\ac{GR}).} AI tasks that use graphs for inference, like common sense reasoning and knowledge-graph question answering. 
    \item \textbf{Graph Abstraction (\ac{GA}).} Transforming graphs into other structures for downstream tasks, for example, generating embeddings with Node2Vec~\cite{Grover2016} or DeepWalk~\cite{Perozzi2014}.
    \item \textbf{Graph Prediction (\ac{GP}).} Learning directly on graphs typically using \ac{GNN}s, including knowledge graph completion and anomaly detection. 
    \item \textbf{Graph Transparency (\ac{GT}).} Uses graphs to improve generative models' transparency, explainability and interpretability. 
\end{itemize}

The taxonomy contextualizes \gai as part of the existing AI landscape and provides a few examples \gai tasks that sit within each \gai category. 
The \gai taxonomy anchors the \gai Benchmark Framework. 
The \gai problem areas guide the development of the benchmark problems within the framework and emphasize the diversity of tasks that the hardware and software systems that run \gai need to support. 

\osullikomment{Add table into Preliminaries}

\section{Challenges in Evaluating \gai}
\label{section:13-graph-based-ai-challenges}


\begin{figure}[h!]
    \begin{center}
        \includesvg[width=\columnwidth]{SVGs/gai_challenges}
        \caption{\gai challenges are a hierarchy of dependencies foundationally limited by graph processing.}
        \label{figure:gai-challenges}
    \end{center}
\end{figure}



A hierarchy of challenges affects \gai (Figure~\ref{figure:gai-challenges}). 
This section briefly explains the core problems at each level of the hierarchy and shows how this thesis's research aim is grounded in addressing these challenges. 

\textbf{Graph Challenges.} All problems in the hierarchy build on the foundational challenge of data dependence in graph processing
Graphs of different shapes and sizes vary the algorithmic complexity and exacerbate the irregular memory accesses and computation patterns that undermine spatial and temporal locality. 
Graphs have no standard accepted representation method, abstractions, or primitive operations.
The resulting landscape consists of competing approaches and no unity of effort in accelerating graph processing generally. 

\textbf{\gai Challenges.} Graph processing is the bottleneck in \gai, meaning that improving \gai performance is directly dependent on elevating the graph processing bottlenecks. 
It is difficult to achieve a universal improvement in \gai because each problem in the taxonomy has different computational characteristics and behaves differently on the same data. 


\textbf{Real-World Challenges.} Graph and \gai processing difficulties are exacerbated by the constraints inherent to deploying them in the real world. 
In particular, \gai solutions must be delivered within the efficiency, effectiveness, and cost constraints. 
For example, where a real-world problem requires an hourly solution but the \gai task that generates it requires ninety minutes to run, compromises must be made to the effectiveness or cost to meet the efficiency constraint. 
Addressing this problem is not as simple as just `accelerating' the \gai task.
For example, hardware accelerators for graph processing are specialized \ac{ASIC} optimized for the specific algorithm, representation, abstraction, and programming model.
Because one system may need to support many \gai tasks, an accelerator will need to improve the performance of all those tasks to be a wise investment. 

\textbf{User Challenges.} The user then must decide whether to trade effectiveness for efficiency or pay a higher cost to achieve the desired results.
Gathering the data to make that decision has historically included choosing between a simple competitive benchmark of systems or an invasive workload analysis to determine which combination of hardware, representation, abstraction model, and algorithm will best meet requirements. 

\textbf{The Evaluation Challenge.}The challenges culminate in the research contribution of this thesis, \textit{determining \textit{if}, \textit{how} and \textit{when} to speed up problem-solving on linked data given a set of operating constraints.}

\section{Overview of the \gai Framework}
\label{section:graph-based-ai-framework}

    
    \begin{figure}[h!]
        \begin{center}
            \includesvg[width=\columnwidth]{SVGs/gai_benchmark_high_level.svg}
            \caption{The high-level view of \gai benchmarking framework}
            \label{figure:gai-framework-overview}
        \end{center}
    \end{figure}



The \gai Benchmark Framework shown in Figure~\ref{Figure:gai-framework-overview} is an Input/Output benchmark that evaluates the performance of a system under test by providing problems, recording the time to solution, and verifying the correctness of an answer. 
The ideal system under test achieves a perfect effectiveness score in minimal time and cost, but all three dimensions will vary in the systems assessed. 

The benchmark is designed so that the system under test is completely independent of the benchmark itself. 
Using the Input/Output paradigm only requires the \ac{SUT} to interact with an \ac{API} (the observation apparatus). It so places no constraints on how a \ac{SUT} conFigures its hardware, software stack, or specific \gai implementation to solve the problem. 

The benchmark definition of the \gai framework is comprised of independent benchmark problem modules. 
Each module contains a collection of datasets designed to stress the efficiency and effectiveness of \ac{SUT} and uses \textit{problem-centric} task definitions to ensure problems are grounded in the real world and to encourage innovative solutions by removing the constraints on kernel-level tasks most existing benchmarks enforce. 

The independence of the benchmark specification from the \ac{SUT} means that an analyst, engineer, or researcher using the benchmark can quickly modify their \ac{SUT} and validate performance changes on a known benchmark. 
The careful design of the component benchmark problems to achieve the `grey-box' analysis of a \gai \ac{SUT} requires a formal methodology for transparency, reproducibility, and verifiability. 


\section{Overview of the \gai Benchmark Problem Development Methodology}
\label{section:graph-based-ai-problem-methodology}
% \begin{landscape}
    
    \begin{figure}[h!]
        \begin{center}
            \includesvg[width=\columnwidth]{Figures/SVGs/benchmark_design_process}
            \caption{\gai Benchmark Design Framework}
            \label{figure:framework}
        \end{center}
    \end{figure}

% \end{landscape}

A formal method for designing benchmark problems ensures the transparency, verifiability, and reproducibility of benchmark problems and guides the extension of the benchmark to include emerging \gai problems. 
Figure~\ref{Figure:framework} shows how the process comprises six steps that generate the components of the benchmark specification and the observation apparatus.

The process begins by identifying a real-world use case to ground the problem in real-world datasets and current approaches. 
A more detailed analysis of the problem yields information about the size and shape of the data and formally specifies the everyday tasks associated with that problem in natural language. 
Formalizing the study's phenomena includes identifying what specific interactions of dataset and task are likely to yield interesting results, providing the first significant contribution to the problem's `grey box' approach. 
Considering the phenomena, exploratory analysis of the domain builds the base of knowledge to make initial hypotheses about behavior. 
The hypotheses drive the data collection to study the phenomena of interest, and deviation from those hypotheses indicates interesting behavior worthy of further study. 
The construction of the apparatus involves collecting and generating the datasets that will be used to trigger the conditions the hypotheses predict will stimulate the phenomena of interest. 
It also involves developing any required \ac{API} extensions to facilitate the problem and evaluate the solutions. 
The benchmark user is responsible for all results analysis beyond the return of the efficiency and effectiveness reports by the benchmark. 

The methodology for designing benchmark problems yields a collection of datasets, tasks, and metrics that will combine to provide a `grey box' view of the \ac{SUT} stimulating phenomena of interest according to the defined hypotheses. 


\section{Thesis Outline}
\label{section:15-Thesis-Outline}
The thesis continues with a review of the \gai taxonomy and key graph and computer architecture concepts in Chapter~\ref{chapter:20-Preliminaries} to identify what makes optimizing \gai approaches hard and why measuring relative performance is difficult. 
Then, Chapter~\ref{chapter:30-Literature-Review} is a detailed review of current \gai benchmarking and workload analysis efforts to characterize the current gaps in benchmark coverage of \gai tasks, datasets, and metrics and identify what causes performance to degrade in \gai systems. 
Chapter~\ref{chapter:40-Graph-AI-Benchmarking} introduces the \gai benchmark frameworks and outlines our methodology for developing \gai problems to constitute the benchmark framework. 
We use the population segmentation problem from the Intelligence and Security domain to demonstrate how our development methodology is applied, describe the implementation of the \gai benchmark prototype, and outline the findings of our initial experiments in chapter~\ref{chapter:50-Community-Detection-Case-Study}.
Finally, chapter~\ref{chapter:60-Conclusion} concludes after identifying specific directions for future work.

%Heilmeier 1, 2, 3, 4

% GraphAI Taxonomy

% GraphAI Challenges

% GraphAI Benchmark Framework 

% GraphAI Development Methodology

% Case Study Use Case

% Population Segmentation Benchmark 

% Experiments 

% Results

% Contributions 

